% Transcription — fonds Alexandre Grothendieck, shelfmark 156-4, batch 3 (pages 41-60).
% Established from the facsimile archives/batches/156-4/batch-03.pdf (a copy of
% 156-4.pdf fetched through the project's relay; no cover sheet in the batch, so
% PDF page k = folder page 40+k), per the skill transcribe-grothendieck. The
% folder is eighty-eight pages in five batches; this is the third. Batches 1
% and 2 were in the repository when this pass was made and their notation is
% followed (batch 2 ends on subdivisions of figures, p. 40, and page 41 opens
% on a Proposition about them).
%
% Pass: Opus 5.5 (claude-opus-5-5), 2026-09-26 — first pass, unchecked against the pages by a human.
%
% Pages skipped: none. Pages summarised: none. All twenty leaves are his
% mathematics, in ink, fast drafting throughout; the lower half of page 45,
% pages 47, 49, 51 and 56 carry dense oblique left-margin notes read only in
% fragments.
%
% His pagination 41-60 runs at the top of the leaves and coincides with the
% archivists'. One date in his hand: « 13 juin » (underlined) opening page 52;
% it is set as a subsection heading, as batch 2 does for « 12 juin », with a
% \note. No other date on these pages. His numbered runs: Cont 7 (p. 41, the
% « Cont 6 » it invokes being batch 2's p. 39), Cor. 1 (p. 41), a cancelled
% Cor. 2 (p. 42), Cont 8 (p. 43), a cancelled « Cont 8 » heading (p. 44); then,
% after « Donc, nouveau départ ! » (p. 45), a new run of axioms C 1 (p. 45),
% C 2 and C 3 (p. 46, both numerals written over others), C 4 (p. 47, with
% Cor. 1 and Cor. 2), Cor. 3 (p. 48), C 5, C 6, C 7 (p. 49); page 52 recalls
% the older Cont 3, 4, 6, 7 as still to be expressed, then opens a further
% framework, « Contrée (quatrième mouture !) », with a third run C 1 a)-b)
% (p. 53), C 2 (p. 54, with Corollaires 1-2), C 3 (p. 58, a first version
% cancelled). Page 51's proof ends « (C'était l'ancien Cont 8) »; page 59's
% « cf. remarque page 58 » and page 60's margin « Cont 2 de p. 33 » are his
% cross-references, given as plain page numbers.
%
% What the batch is about. Continuing the first draft of chapter IV. Pages
% 41-44: in batch 2's framework (multistrates M, relation ≪, lieux, ombres X̃
% and open ombres X̃°), a Proposition characterising a subdivision F' of F
% among refinements (a ⟺ no X ∈ M refining F and disjoint from F' ⟺ no lieu
% of the ombre F̃ disjoint from F'), needing a new axiom Cont 7 (distinct
% compatible X, Y have disjoint open ombres, lieu by lieu); an attempt to
% map F' → F by open strata, which stalls (« il semblerait qu'il manque
% quelque chose »); Cont 8, the refinement R_{X'} induced by Y on X' as the
% inf of F_Y and F_{X'}, with Ỹ ∩ X̃' = ∪ Ỹ', and a Cor. Ỹ° ⊂ X̃°; a Corollary
% that each open stratum of a refinement lies in a unique open stratum of F;
% then (p. 44) the observation, on pseudo-discs, that two disjoint interior
% lieux would not define a subdivision, so that refinement cannot be the
% primitive notion. Pages 45-51, « nouveau départ » : a contrée is a triple
% (M, 𝔉, Σ) with Σ the relation « F' subdivise F »; C 1 (elementary figures
% F_X, order X ≤ Y ⟺ F_X ⊂ F_Y), figures are closed parts, compatibility,
% C 2 (compatibility tested on strata), C 3 (closed parts of a figure are
% sub-figures), empty figure and empty contrée (p. 46); refinement G ≪ F :=
% G ⊂ a subdivision of F, C 4 (a subdivision of a sub-figure extends as
% G' ∪ (F ∖ G), with G' ∩ F ⊂ G), ≪ an order, Cor. 3 (G ⊂ F ⟺ G compatible
% with F and G ≪ F), disjoint figures (pp. 47-48); C 5 (disjointness
% unchanged by subdivision), C 6 (induced subdivisions F'|G), C 7 (gluing
% subdivisions), and the Proposition and Corollary that the induced
% refinement G' = {X ∈ F' | X ≪ G} is the inf of F' and G, proved with four
% diagrams of subdivisions (pp. 49-51). Pages 52-60, 13 June: a « quatrième
% mouture » taking the set 𝔉 of figures as primitive, with two orders ≤
% (sub-figure) and ≼ (subdivision): C 1 (non-empty; bounded families have a
% Sup), the empty figure, irreducible figures = multistrates M, strates F̄
% (p. 53); C 2 (F = Sup of its strata), 𝔉 ≅ a set Φ ⊂ 𝔓(M) of parts with
% conditions a)-c), and conversely (pp. 54-55); sub-figures as closed parts,
% F ∧ G, lieux as minimal non-empty figures, disjoint figures and the
% bijection 𝔉_{≤F} × 𝔉_{≤G} ≅ 𝔉_{≤F∪G} (pp. 56-57); C 3 (compatibility
% tested on strata, p. 58); irreducible components, figures finite and of
% finite type, and the reconstruction of (𝔉, ≤) from (M, ≤, R) with R a
% reflexive symmetric compatibility relation (pp. 59-60).
%
% Notation fixed once. As in batches 1-2: \mathcal{M}, \mathcal{L}, \mathfrak{F},
% \mathfrak{P}; ombres \widetilde{X}, open ombres \widetilde{X}^{\circ}; refinement
% ≪ (\ll), the triple chevron between figures \lll where he writes it (p. 41);
% his barred cursive capital (a part of M, and on p. 51 an auxiliary figure)
% \Phi, as in batch 1; the curly ≼ of subdivision (p. 52 on) \preccurlyeq;
% the set of strata of F, his barred F, \overline{F}; on pp. 50-51 the
% barred and tilded figures of the proof \overline{F}', \widetilde{F},
% \vec{F} as drawn. Compatibility is R (p. 59; \mathcal{R} in the
% cancelled proposition of p. 47). His underlining is set \emph.
%
% Hard pages: 44 (the pseudo-disc paragraph), 45-47 (margins), 51 (four
% diagrams, described in one note, not redrawn), 59 (connective prose).
% Apparatus: 133 \ill{}, 27 \uncertain{}. Variances on the page, not
% repaired and flagged: F ∪ F_X for F' ∪ F_X (p. 41); « qui raffine G' » for
% G (p. 50); G' ≪ F for G' ≪ G (p. 49); H ∩ F for H' ∩ F (p. 57); A ∈ M for
% A ⊂ M (p. 60). Readings flagged and not settled include « quatrième » in
% the heading of p. 52, the direction of « F ≪ G » in the margin of p. 47,
% the insertion on 𝔓_f(𝔉) (p. 57), and « l'ancien Cont 8 » (p. 51).

\documentclass[11pt,a4paper]{article}
\input{../preamble/grothendieck}

\folder{156-4}
\batch{3}
\pages{41}{60}
\dating{1986}
\watermark{Édition de démonstration}
\foldertitle{[Chapitre] IV. Analysis situs (première mouture) : notes manuscrites (10/06/1986)}

\begin{document}

\page{41}
\emph{Prop.} Soit $F'$ un raffinement de $F$. Conditions \emph{équivalentes} :
\begin{enumerate}
  \item[a)] $F'$ est une subdivision
  \item[b)] Il n'existe pas de $X \in \mathcal{M}$ qui raffine $F$, et disjoint de $F'$
  \item[c)] \struck{Il n'existe pas \ill{}} \struck{$\widetilde{F}' = \widetilde{F}$, i.e., \ill{} il n'existe pas de lieu dans $F$} Il n'existe pas de lieu dans l'ombre $\widetilde{F}$, qui soit disjoint de $F'$ \marginal{(avec Cont 6)} i.e. $\forall x \in \widetilde{F} \smallsetminus \widetilde{F}'$, $\exists\, y \in \widetilde{F}'$ tel que $x$ et $y$ \emph{non compatibles}, ou (positivement) : si $x \in \widetilde{F}$ est compatible à \struck{\ill{}} \uncertain{tout} élément de $\widetilde{F}'$, il est dans $\widetilde{F}'$.
\end{enumerate}
\note{la condition c) est d'abord commencée puis barrée sur deux lignes ; « (avec Cont 6) » est ajouté au-dessus de la ligne, en fin de ligne}

\emph{Dém.} a $\Rightarrow$ b, car si $X \in \mathcal{M}$ raffine $F$ et est disjoint de $F'$, alors la figure $F \cup F_X$ est un raffinement de $F$, strictement plus grand que $F'$.
\note{on attend $F' \cup F_X$ ; la page porte $F \cup F_X$}

b) $\Rightarrow$ c) tautologique

c) $\Rightarrow$ a) Soit \ill{}, $F' \subset F'' \lll F$, et soit $X \in F''$, tel que $X \notin F'$, et soit $x \in \widetilde{X}^{\circ}$. Alors $x \in \widetilde{F} \smallsetminus \widetilde{F}'$, et \struck{\ill{}} je dis que $x$ est compatible à $F'$ donc c) n'est pas vérifié.
\marginal{à vérifier}
\note{« à vérifier » est écrit en marge gauche, avec une parenthèse qui embrasse les deux dernières lignes}

Mais on a utilisé que $x$ est compatible à $F'$, il faut un axiome pour l'assurer.

\emph{Cont 7} \quad Soient $X, Y \in \mathcal{M}$, \add{$X$ et $Y$ compatibles (donc $\widetilde{X}^{\circ} \cap \widetilde{Y}^{\circ} = \emptyset$)}, $X \neq Y$. Alors tout $x \in \widetilde{X}^{\circ}$ est disjoint de tout $y \in \widetilde{Y}^{\circ}$ (ou, ce qui revient au même par Cont 6, \struck{\ill{}} tout $x \in X^{\circ}$ est disjoint de $Y$)
\note{triple trait vertical en marge de Cont 7 ; l'insertion est écrite au-dessus de la ligne et reliée au texte par un trait}

\emph{Cor. 1} \quad Soient $F$ une figure, $F'$ une sous-figure et $X \in F \smallsetminus F'$. Alors $\forall x \in X^{\circ}$ est \add{disjoint} \struck{compatible} de $F'$.
\note{« disjoint » est écrit au-dessus de « compatible » biffé ; lecture douteuse}

\page{42}
\struck{\emph{Cor. 2} Soit $F$ une figure. Deux lieux dans l'ombre de $F$, qui appartiennent à des ombres ouvertes de strates distinctes, sont disjoints.}
\note{deux lignes barrées en tête de page, qui recommencent un « Cor. 2 » ; le numéro « 42 » est écrit au milieu de la première}

Soient $F$ une figure, \struck{$X \in \mathcal{M}$ un} $F'$ un raffinement \struck{\uncertain{raffinement de $F$}} de $F$. Je veux définir une application
\[
  F' \to F
\]
en associant : $\forall X' \in F'$ (l'unique) $X \in F$ tel que
\[
  \widetilde{X}'^{\circ} \subset \widetilde{X}^{\circ}
\]
(ou encore le plus petit $X \in F$ tel que $\widetilde{X}'^{\circ} \subset \widetilde{X}$). Savons-nous déjà qu'un tel $X$ existe, i.e. que la relation \uncertain{d'équivalence} \ill{} sur l'ombre de $F'$, est plus fine que celle induite par l'ombre de $F$ ? Soit $X$ un élément de $F$ \emph{minimal}, parmi ceux satisfaisant
\[
  X' \ll X
\]
(existe, à cause de l'axiome des \uncertain{lieux}), savons-nous \uncertain{dès lors} $\widetilde{X}'^{\circ} \subset \widetilde{X}^{\circ}$ ? \struck{Évid} On sait que $\widetilde{X}' \subset \widetilde{X}$ donc si $x \in \widetilde{X}'^{\circ} \subset \widetilde{X}'$, i.e. $x \in \widetilde{Y}^{\circ} \subset \widetilde{Y}$ pour $Y \leq X$, on veut prouver que $Y = X$. Il semblerait qu'il manque quelque chose.
\note{« i.e. » sous $\widetilde{X}'$ est ajouté sous la ligne ; le $Y$ de $\widetilde{Y}^{\circ}$ est écrit sur une autre lettre}

\page{43}
\marginal{\uncertain{Relation} avec \uncertain{raffinement induit}}
\emph{Cont 8} \quad Soient $X$, $Y$, $X'$ dans $\mathcal{M}$, $Y \ll X$, $X' < X$. Considérons l'ens. des $Y' \in F_Y$ (i.e. $Y' \leq Y$) tels que $Y' \ll X'$ :

\[
  \begin{array}{ccc}
    Y & \ll & X \\
    \vee & & \vee \\
    Y' & \ll & X'
  \end{array}
\]
\note{les deux signes verticaux sont des $\leq$ couchés : $Y' \leq Y$, $X' \leq X$}

(lesquels forment une sous-figure $R_{X'}$ \ill{} multistrate $\in F_Y$, \struck{qu'on appelle} raffinant $X'$, et qu'on appellera le raffinement de $X'$ induit par $Y$ sur $X'$). Alors tout raffinement commun $Z$ de $Y$ et de $X'$ raffine un des $Y'$, i.e. raffine $R_{X'}$. (Donc toute figure qui raffine à la fois $F_Y$ et $F_{X'}$, raffine $R_{X'}$, qui est donc \struck{la} \struck{bonne} la borne inf. de $F_Y$ et $F_{X'}$ pour la relation $\ll$) En particulier,
\[
  \widetilde{Y} \cap \widetilde{X}' = \bigcup_{Y' \in R_{X'}} \widetilde{Y}'
\]
\note{la marge gauche porte, en oblique, trois mots de lecture douteuse ; la lettre $R$ de $R_{X'}$ est barrée d'un trait, comme dans tout ce passage}

\emph{Cor} \quad \struck{\ill{}}, Supposons que $\nexists\, X_1 \in \mathcal{M}$, \quad $Y \ll X_1 < X$ (c'est la situation de \ill{}). Alors $\widetilde{Y}^{\circ} \subset \widetilde{X}^{\circ}$.

En effet, si $x \in \widetilde{Y}^{\circ}$, s'il avait $x \in \widetilde{X}'$ pour $X' < X$, on aurait $x \in \widetilde{Y}'$ avec $Y' \in R_{X'}$, donc $Y' = Y$ \struck{$X$} \uncertain{par déf.} de $\widetilde{Y}^{\circ}$, donc $Y \ll X' < X$, contradiction.

\page{44}
\emph{Corollaire} \quad Si \struck{\ill{}} $Y \ll X$, alors \struck{la relation} toute ombre ouverte d'une \struck{\ill{}} \add{strate \ill{} $Y$} est contenue dans l'ombre ouverte d'une strate de $X$. Plus généralement et \uncertain{précisé}, pour \struck{\ill{}} figure $F$ et \add{\ill{}} raffinement $F'$ de $F$, \struck{\ill{}} \struck{\ill{} de l'ombre ouverte d'une} strate ouverte de $F'$ est contenue dans une (unique) strate ouverte de $F$.
\note{la phrase a été corrigée en cours d'écriture ; seule la ligne principale est lue}

\struck{\emph{Cont 8 \ill{}}}
\note{un titre biffé, suivi d'un grand trait oblique}

Je viens de m'apercevoir que dans le cas des pseudo-disques, \struck{\ill{}} deux lieux disjoints \add{$x$, $y$} de l'intérieur (lesquels définissent deux \uncertain{deux} « raffinements » \uncertain{induits} pseudo-disques), qu'il existe une « subdivision » des disques \ill{} ils soient des \add{de $F$} \emph{noeuds}. [\emph{NB.} Les \emph{noeuds} d'une figure \add{$F$} sont les multistrates \add{de $F$} réduites à un lieu, i.e. les éléments de $F \cap \mathcal{L}$] Donc \struck{\ill{}} il n'y a pas lieu de considérer $\{x, y\}$ comme un \ill{} « raffinement », ou qu'il ne définit pas une \emph{subdivision}. Donc il est inapproprié
\note{« de $F$ » est ajouté en marge gauche devant « réduites » ; au-dessus de « des » un second « de $F$ » ; la construction de la première phrase est celle de la page}

\page{45}
de définir une figure $F'$ comme \emph{raffinement} de $F$, si chaque $X' \in F'$ est un raffinement d'un $X \in F$. Donc il vaut mieux prendre comme notion primitive celle de \emph{subdivision} d'une \emph{figure}, (qui \ill{} ; celle d'une multistrate), et définir un raffinement d'une figure comme une sous-figure d'une subdivision.

\section{Donc, nouveau départ !}

Une contrée $C$ est définie par un triple
\[
  (\mathcal{M}, \mathfrak{F}, \Sigma)
\]
où $\mathcal{M}$ est un ens. (l'ens. des \emph{multistrates}), $\mathfrak{F}$ un ensemble de parties de $\mathcal{M}$ \add{dites \emph{figures}}, ($\mathfrak{F} \subset \mathfrak{P}(\mathcal{M})$), et $\Sigma$ une relation sur $\mathfrak{F}$ dite : $F'$ \emph{subdivise} $F$. Axiomes \add{\ill{} des éléments de $\mathcal{M}$}
\marginal{Si $F \in \mathfrak{F}$ les $X \in F$ sont \ill{} strates de $F$}

\emph{C 1} \quad Pour tout $X \in \mathcal{M}$, $\exists$ une \add{\ill{}} plus petite figure $F_X$ telle que $X \in F$. \struck{\ill{}} figure \struck{qui} \add{i.e.} \ill{} \emph{figure élémentaire} définie \uncertain{par} $X$ \add{si $X \neq Y$, $F_X \neq F_Y$}. On écrit
\[
  X \leq Y \overset{\mathrm{def}}{\Longleftrightarrow} F_X \subset F_Y \quad \text{i.e. } X \in F_Y
\]
et on trouve une relation d'ordre* $\leq$ sur $\mathcal{M}$.
\marginal{* On veut que ce soit une relation d'ordre \ill{}}
\marginal{C'1 \quad $\exists\, F \in \mathfrak{F}$ \ill{}}
\note{C 1 est marqué d'un double trait vertical en marge ; une insertion en petit, entre les lignes, porte « $F_X \neq F_Y$ » et « $X \neq Y$ » et n'est lue qu'en partie ; en marge gauche, en bas, un « C'1 » précédé de deux traits, dont le texte n'est lu qu'en fragments}

\page{46}
\struck{\ill{} $F_X$}

\ill{} $F$ est une figure, alors $X \in F$ équivaut à : $F_X \subset F$, qui implique que pour tout $Y \in F_X$ (i.e. $Y \leq X$) on a $F_Y \subset F$ i.e. $Y \in F$. Ainsi, $F$ est une partie de $\mathcal{M}$ \emph{fermée} pour $\leq$.
\note{la marge gauche, en haut, porte une longue note oblique de six ou sept lignes, en partie biffée (elle commence par « Donc, $\mathcal{M}$ s'identifie \ill{} », et mentionne $F_X$ et « relation d'ordre ») ; elle n'est lue qu'en fragments}

On dit que deux figures \add{$F$, $G$} sont \emph{compatibles} si $F \cup G$ est une figure, et que $X, Y \in \mathcal{M}$ sont \emph{compatibles} si $F_X$, $F_Y$ le sont.

\emph{C 2} \quad Pour que $F$ et $G$ soient compatibles, il f. et s. que $\forall X \in F$, $Y \in G$, $X$ et $Y$ le soient.
\marginal{correspond à Fig 4}
\note{double trait vertical en marge ; le chiffre 2 est écrit sur un autre}

\struck{\emph{Cor.} \ill{} pour que $X, Y \in \mathcal{M}$ soient compatibles, il f. et s. qu'ils \ill{} appartiennent à une figure}

\struck{\emph{C 3} \ill{}}

\emph{C 3} \quad Soit $F$ une figure, alors toute partie de $F$ fermée pour $\leq$ est une figure (i.e. pour qu'une partie $F'$ de $F$ soit une figure, il f. et il suffit qu'elle soit fermée pour $\leq$).
\marginal{\uncertain{intersection}}
\note{double trait vertical en marge ; le chiffre 3 est écrit sur un autre ; la note marginale est reliée par une flèche à C 3, une autre flèche remonte vers C 2}

(\emph{Sous-figures}) Si $\mathcal{M} \neq \emptyset$ cela implique que la partie vide de $\mathcal{M}$ est une figure. Si $\mathcal{M} = \emptyset$, \struck{cela implique que} on aura $\mathfrak{F} = \emptyset$ \struck{\ill{}} ou $\mathfrak{F} = \{\emptyset\}$, et on suppose que c'est le deuxième cas : \emph{figure vide existe toujours}
\marginal{\emph{Contrée vide} si $\mathcal{M} = \emptyset$}

\page{47}
\emph{Définition} \quad Soient $F$, $G$ deux figures. On dit que $G$ est un \emph{raffinement} de $F$, s'il existe une \struck{raffinement} subdivision $F'$ de $F$, telle que $G \subset F'$ \add{sous-figure}.
\marginal{On écrit $F \ll G$}
\note{le sens du signe de la note marginale est douteux ; aux pages 48 et 49 il écrit $G \ll F$ pour « $G$ raffine $F$ »}

On veut que ce soit une relation d'ordre, et pour ceci on va poser

\emph{C 4} \quad Soient $F$ une figure, $G$ une sous-figure, \add{(i.e. $G$ figure et $G \subset F$)} $G'$ une subdivision de $G$. Alors \struck{il existe une subdivision $F'$ de $F$,} \add{$G' \cup (F \smallsetminus G)$ est une} subdivision \add{de $F$}, et la réunion est disjointe (i.e. \struck{\ill{}}
\[
  \boxed{G' \cap F \subset G}
\]
donc $G' \cap F = G' \cap G$)
\marginal{Corriger : $\exists$ subdivision qui induit $G'$ + la relation $G' \cap F \subset G$}
\marginal{OK}
\note{double trait vertical en marge de C 4 ; « $G' \cup (F \smallsetminus G)$ est une » est écrit en gros caractères au-dessus de la ligne biffée ; en marge gauche une note de cinq lignes, reliée par une flèche à la formule encadrée (« \ill{} multistrates \ill{} qui sont dans $G$ »), n'est lue qu'en fragments ; la correction de droite est écrite en petit sous « donc »}

\emph{Cor. 1} \quad Toute \struck{raffinement} subdivision d'une sous-figure de $F$ est induite par une « subdivision » de $F$.
\marginal{OK}

\emph{Cor. 2} \quad La relation de raffinement entre figures
\[
  F \ll G
\]
est une relation d'ordre.
\marginal{OK.}

\struck{\emph{Proposition} (\ill{}) Soient $X \in \mathcal{M}$, $G \in \mathfrak{F}$. Pour que $X \in G$ \struck{$F_X \subset F$}, il f. et s. que $X$ et $G$ compatibles (i.e. $F = F_X \cup G$ une figure) et $X \ll G$ i.e. $F_X$ raffine $G$. \struck{\ill{}} Supposons en effet que $X$ et $F$ compatibles, et que $F_X$ raffine la sous-figure $F_Y = \ill{}$. Donc $\exists$ subd. $G'$ de}
\marginal{C'est la \uncertain{définition} \ill{} en $\ll$ et $\mathcal{R}$}
\note{toute cette proposition, qui déborde sur la page 48, est barrée de grands traits obliques ; la lettre $G$ y est plusieurs fois écrite sur $F$}

\page{48}
\struck{$G$ telle que $F_{X'} \subset G'$ i.e. $X' \in G'$. Cette subdivision est induite par subdivision $G' \cup (F \smallsetminus G)$, qui \ill{} disjointes, et comme $X' \in G'$ \ill{}}
\note{fin de la proposition barrée de la page 47, barrée elle aussi de traits obliques et entourée}

\emph{\struck{Proposition} Cor. 3} \quad Soient $G$, $F$ deux figures. Pour que $G \subset F$ \add{(i.e. $G$ sous-figure de $F$)}, il f. et s. que $G$ \emph{compatible} à $F$ (i.e. $F \cup G$ une figure, soit $H$) et $G \ll F$, i.e. $G \subset F'$, $F'$ subd. de $F$.

En effet, par \struck{\ill{}} C 4, \struck{\ill{}}
\[
  G \subset F' \cap H \overset{\text{C 4}}{\subset} F
\]
donc $G \subset F$.
\[
  \begin{array}{c}
    F \subset H = F \cup G \\
    \uparrow{\scriptstyle \text{subd.}} \\
    F' \\
    \cup \\
    G
  \end{array}
\]
\note{le petit schéma est à gauche de la formule ; après $F' \cap H$ quelques signes surchargés sont biffés}

En particulier, on trouve : si $X \in \mathcal{M}$, $F \in \mathfrak{F}$, alors
\[
  X \in F \Longleftrightarrow
  \begin{cases}
    X \text{ comp. à } F \\
    X \text{ raffine } F
  \end{cases}
\]
donc si $X, Y \in \mathcal{M}$, alors
\[
  X \leq Y \Longleftrightarrow
  \begin{cases}
    X \text{ comp. à } Y \\
    X \ll Y
  \end{cases}
\]
\marginal{\ill{} de C 4}
\note{dans la seconde accolade, « raffine » est biffé et remplacé par $\ll$}

\emph{Déf} \quad $F, G \in \mathfrak{F}$ sont dites \emph{disjointes} si elles sont compatibles (i.e. $F \cup G \in \mathfrak{F}$) et $F \cap G = \emptyset$. On dit que $X$, $Y$ sont disjoints si $F_X$ et $F_Y$ le sont i.e. $X$ et $Y$ compatibles, et $\nexists\, Z$ avec $Z \leq X, Y$. Si $F$ et $G$ sont disjointes, alors pour $F' \subset F$, $G' \subset G$ sous-figures, elles sont disjointes. De plus ;
\marginal{\emph{NB} $F \cap G$ disjoints $\Leftrightarrow$ $\forall X \in F$, $Y \in G$, $X$ et $Y$ disj. (i.e. \uncertain{compatibles} et distincts)}
\note{la note marginale est écrite en oblique dans l'angle inférieur gauche ; au-dessus de « Déf », une insertion oblique (« $X$ et $Y$ disj. \ill{} distincts ») n'est lue qu'en partie}

\page{49}
\struck{\emph{Cont 5} \quad Si $F$ et $G$ sont deux figures disjointes, et $F'$ \add{de $F$} \struck{\ill{}} subdivisions, alors $F'$ et $G$ sont disjointes. [\emph{NB} $F \cap G$ \ill{} dans $F$ et $G$ compatibles \ill{}]}

\struck{\emph{Cor} Si $F'$ raffine $F$, $G'$ raffine $G$, \ill{} $F$ et $G$ disjointes alors $F'$ et $G'$ disjoints.}
\note{ces deux énoncés sont barrés de grands traits obliques}

\emph{C 5} \quad Soient $F$ une figure, $F'$ une subdivision de $F$. Alors pour toute figure $G$, $G$ est disjointe de $F$ ssi elle est disjointe de $F'$ ;
\note{double trait vertical en marge ; « C » est écrit sur un autre signe ; « $F$ ssi elle est disjointe de $F'$ » est écrit au bout de la ligne, au-dessus de « $G$ est disjointe ; »}

\emph{NB} Il suffit de le prouver pour $G$ élémentaire.

Il résulte alors de \struck{Cont 4} C 4 que $F' \cup G$ est une subdivision de $F \cup G$.

\emph{Cor} \quad Soient $F$ et $G$ figures disjointes, alors pour $F' \ll F$, $G' \ll F$, $F'$ et $G'$ sont disjointes, et $F' \cup G' \ll F \cup G$. Si $F'$ et $G'$ sont des subd. de $F$, $G$, alors $F' \cup G'$ en est une de $F \cup G$.
\note{on attend $G' \ll G$ ; la page porte $G' \ll F$}

\emph{C 6} (\emph{Subdivisions induites}) \quad Soient $F$ une figure, $F'$ une subdivision de $F$, $G$ une sous-figure de $F$. Alors il existe une unique subdivision $G'$ de $G$, telle que
\[
  G' \subset F'
\]
(subdivision induite) \add{(et $G'$ est formé des $X \in F'$ tels que $X \ll G$)}.
\note{double trait vertical en marge}

\emph{NB} Transitivité de l'induction. \emph{Notation} $F'|G$
\marginal{\uncertain{Doublé}}

\emph{C 7} \quad (\struck{\ill{}} \emph{Recollement} de subdivisions) Soient $F$ une figure, \struck{\ill{}} \add{$F_i$ des sous-figures, \ill{}} tels que $F = \bigcup F_i$. Soit, pour tout $i$, $F'_i$ une subdivision de $F_i$, et supposons
\note{double trait vertical en marge ; le titre entre parenthèses est écrit au-dessus de la ligne}

\page{50}
que pour \struck{\ill{}} $i, j \in I$, $i \neq j$,
\[
  F'_i \,|\, F_i \cap F_j = F'_j \,|\, F_i \cap F_j
\]
Alors il existe une unique subdivision $F'$ de $F$ telle que $F' \,|\, F_i = F'_i$ \quad $\forall i \in I$.

\emph{Cor} \quad Pour que deux subdivisions \add{$F'$, $F''$ de $F$} soient égales, il f. et s. qu'elles induisent la même subdivision sur chaque sous-figure élémentaire $F_X$, i.e. sur chaque multistrate.

\emph{\struck{Cor} Proposition} \quad Soient $F$ une figure, $G$ une sous-figure, $F'$ un raffinement de $F$, et
\[
  G' = \lbrace X \in F' \mid X \ll G \rbrace
\]
Je dis que c'est donc une sous-figure \add{\ill{}} de $F'$, et elle raffine $G$ et c'est la plus grande sous-figure de $F'$ qui raffine $G'$.
\marginal{\emph{NB} \ill{} $F' \in \ill{}$ \ill{} $G'$ \ill{} raffinement induit}
\note{on attend « qui raffine $G$ » ; la page porte $G'$. En marge gauche, sous la note oblique, un petit schéma : $F' \ll F$ au-dessus de $G' \ll G$, reliés par deux signes d'inclusion verticaux}

\struck{Cela dit,} Il y a qqch à prouver, car le seul fait que les $X \in G'$ raffinent $G$ n'implique pas que $G'$ raffine $G$. Mais on a une \struck{raffinement} \add{subdiv.} $\overline{F}'$ de $F$ tel que
\[
  \begin{array}{ccc}
    F' \subset \overline{F}' & \xrightarrow[s]{\sim} & F \\
    \cup & & \cup \\
    \overline{G}' & \xrightarrow[s]{\sim} & G
  \end{array}
\]
\marginal{considérons la subd.}
(C, 6) induit \struck{\ill{}} $G$, $\overline{G}'$, formé des $X \in \overline{F}'$ tels que $X \ll G$. \struck{\ill{}} Par définition, on a donc
$G' = F' \cap \overline{G}'$, qui est \struck{\ill{}} une sous-figure de $\overline{G}'$ donc raffine $G$.
\note{dans le schéma, la flèche marquée $\sim$ et $s$ figure une subdivision ; les inclusions verticales sont dessinées comme des $\cup$}

\page{51}
\emph{Corollaire} \quad Sous les conditions précédentes, $G'$ est la Inf de $F'$ et $G$, dans l'ens. des figures pour la \add{relation} \struck{\ill{}} de raffinement. En d'autres termes, toute figure \add{$\Phi$} qui raffine à la fois $F'$ et $G$, raffine $G'$.
\note{trait vertical en marge de l'énoncé ; la lettre $\Phi$, ajoutée au-dessus de « figure », rend sa capitale cursive barrée, comme au premier lot}

On utilise que $F' \subset \overline{F}'$, où $\overline{F}'$ subd. de $F$, et que $\Phi \subset \overline{F}''$, où $\overline{F}''$ subd. de $F'$, \struck{\ill{}} On va \ill{} \add{(d'où $\Phi \ll G'$)} $\Phi \subset \overline{G}'$. Soit $X \in \Phi$, on a $X \in \overline{F}'$, \struck{et $X$} il faut prouver $X \in \overline{G}'$, i.e. \ill{} pour \uncertain{tt} élément de $\overline{F}'$ \emph{qui raffine $G$}, est dans $\overline{G}'$. Mais $\overline{F}'$ \emph{est induit par une subd.} $\widetilde{F}$ de $F$ (C 4 revu et corrigé -- cf. ci-con.) donc $X \in \widetilde{F}$ et $X$ raffine $G$, donc \struck{\ill{}} $X \in \widetilde{G}$, \struck{subd.} subd. de $G$ induite par la subd. $\widetilde{F}$ de $F$, donc
\[
  X \in \overline{F}' \cap \widetilde{G} = \overline{G}' , \quad \text{qed,}
\]
ouf ! (C'était \uncertain{l'ancien} Cont 8)
\note{la démonstration est écrite dans la moitié droite de la page, la moitié gauche portant quatre schémas successifs. Le premier (\ill{} $\to F' \subset \overline{F}' \simeq F$ au-dessus de $\overline{G}' \simeq \overline{G}' \subset \overline{G} \simeq G$, avec inclusions verticales) est surchargé et barré de traits obliques. Le deuxième aligne $\Phi \subset \overline{F}'' \simeq F' \subset \overline{F}' \simeq F$ au-dessus de $\overline{G}'' \simeq G' \subset \overline{G} \simeq G$, flèches verticales descendantes. Le troisième, $\overline{F}' \subset \widetilde{F} \simeq \vec{F} \simeq F$ au-dessus de $\widetilde{G} \simeq \overline{G} \simeq G$. Le quatrième, entouré d'un grand arc, est un cube : $\Phi \subset \overline{F}'$, $\widetilde{F}$, $\vec{F}$, $F$ en haut, $F'$ au milieu, $\overline{G}'$, $\widetilde{G}$, $\overline{G}$, $G$, $G'$ en bas, reliés par des flèches marquées $\simeq$ (subdivisions) et des flèches verticales ; il n'est pas redessiné. Le « ci-con. » renvoie sans doute à la correction de C 4 en marge de la page 47}

\page{52}
\subsection{13 juin}

\note{la date, soulignée, est de sa main en tête de la page, en début de ligne}
On n'a pas encore exprimé tous les axiomes \emph{Cont.} 1, \ldots, il reste \uncertain{pour compte}
\begin{itemize}
  \item[\emph{Cont 3}] dim finie
  \item[\emph{Cont 4}] « admissibilité » -- mais résulte déjà \ill{} \struck{des} axiomes C.
  \item[\emph{Cont 6}] \add{\uncertain{critères} de} disjonction via les lieux -- on y reviendra
  \item[\emph{Cont 7}] \emph{Deux lieux de strates ouvertes différentes sont disjoints.}
\end{itemize}

Mais avant de continuer, je vais encore une fois changer d'optique ! Prendre comme ens. de base $\mathfrak{F}$, l'ens. des figures, avec les deux relations d'ordre $\leq$ (ou l'inclusion) et $\preccurlyeq$ (la subdivision)

\section{Contrée (\uncertain{quatrième} mouture !)}

\note{titre de sa main au milieu de la page, « Contrée » souligné}

C'est un ens. $\mathfrak{F}$ d'objets, appelés \emph{figures}, et deux relations d'ordre $\leq$ et $\preccurlyeq$ sur $\mathfrak{F}$ :
\[
  \mathcal{C} = (\mathfrak{F}, \leq, \preccurlyeq)
\]
satisfaisant à des axiomes qu'on va énumérer. Si $F$ et $G$ sont deux figures et $F \leq G$, on dit que $F$ est une \emph{sous-figure}, et on écrira aussi souvent $F \subset G$. Si $F \preccurlyeq G$, on dit que $F$ est une \emph{subdivision} de $G$, ou que $F$ subdivise $G$.

\page{53}
D'abord, on va dégager les axiomes pour $\leq$

\emph{C 1} \quad a) $\mathfrak{F} \neq \emptyset$

\struck{\ill{}} b) Si $F \in \mathfrak{F}$, et $(F_i)_{i \in I}$ est une famille d'objets de $\mathfrak{F}$ avec $F_i \leq F$, alors $\operatorname{Sup}_i F_i$ existe. (Toute famille majorée de figures a une borne sup.)
\note{double trait vertical en marge de C 1 ; un signe $\leq$ est écrit en indice sous le $\mathfrak{F}$ de « $F \in \mathfrak{F}$ »}

\emph{Corollaire} \quad $\mathfrak{F}$ a un plus petit élément (savoir, la Sup de la famille vide d'objets)

\emph{NB} \quad En présence de \struck{\ill{}} C 1 b), la validité de a) \emph{équivaut} à celle du corollaire.

Le plus petit élément est appelé \emph{figure vide}, noté $\emptyset_{\mathfrak{F}}$ ou simplement $\emptyset$, \ill{} si confusion n'est à craindre

\struck{\ill{} des figures compat. \emph{Définition} \emph{figures irréductibles} $F$ : telles que si $F = \operatorname{Sup}$}
\note{ces lignes sont encadrées et barrées de traits obliques}

Figures \emph{compatibles}, ens. de figures \emph{compatibles} \add{deux à deux} \struck{que} (ou \emph{configurations}) : qui est majoré pour $\leq$. Figure engendrée, ou \emph{réunion} ($F \cup G$, $\bigvee_i F_i$, $\operatorname{Sup} F_i$, parfois noté $F \cup G$ -- $\bigcup_i F_i$ cf plus bas).
\note{en fin de ligne, au-dessus de « Figures », un mot biffé}

\emph{Déf} \quad \emph{Figure irréductible} $F$ : telle que
\[
  F = \operatorname{Sup}_{i \in I} F_i \Longrightarrow \exists\, i \in I, \ F_i = F .
\]

\struck{\emph{C 2}} Les figures irréductibles sont \add{dites} des \emph{multistrates}. \add{L'ens. des multistrates est noté $\mathcal{M}$}, \struck{les} figures irréductibles

\page{54}
majorées par $F$ s'appellent les \emph{strates} de $F$. (Donc, \struck{\ill{}} strate de $F$ est une multistrate) L'ens. des strates de $F$ \struck{s'appellera} sera noté $\overline{F}$ :
\[
  \struck{\ill{}} \quad \mathcal{M} \subset \mathfrak{F}
\]
\[
  F \mapsto \overline{F} \qquad \mathfrak{F} \to \mathfrak{P}(\mathcal{M}) .
\]

\emph{C 2} \quad Toute figure est la borne supérieure de ses strates :
\[
  F = \operatorname*{Sup}_{X \in \overline{F}} X
\]
\note{double trait vertical en marge}

\emph{Corollaire 1} \quad L'application $F \mapsto \overline{F}$
\[
  \mathfrak{F} \to \mathfrak{P}(\mathcal{M})
\]
est injective, et induit un iso \ill{} d'ens. ordonnés de $\mathfrak{F}$ avec l'ens. $\Phi \subset \mathfrak{P}(\mathcal{M})$ formé des $\overline{F}$. Cet iso transforme $\operatorname{Sup}_i$ en $\bigcup_i$.

\emph{Corollaire 2} \quad L'ens. $\Phi$ de parties de $\mathcal{M}$ satisfait les deux conditions,
\begin{enumerate}
  \item[\struck{a)}] \struck{$\Phi \neq \emptyset$}
  \item[a)] Si $A \in \Phi$, et $(A_i)_{i \in I}$ est une famille d'éléments de $\Phi$, avec $A_i \subset A$, alors $\bigcup_i A_i \in \Phi$.
  \item[b)] $\forall X \in \mathcal{M}$, il existe un plus petit $A = A_X$ tel que $X \in A$ (a fortiori, $\Phi$ recouvre $\mathcal{M}$).
\end{enumerate}
\marginal{\uncertain{Ceci}}
\note{trait vertical en marge du Corollaire 2 ; le « a) » est écrit sur un « b) »}

\page{55}
\struck{c) $A \in \Phi$ \ill{} [car \ill{} vu b), c'est équivalent à $\mathcal{M} \neq \emptyset$]}
\struck{\emph{Prop.} équivalente : \ill{} $\mathcal{M} \neq \emptyset$ \ill{} $\emptyset \in \Phi$ \ill{}}
\note{deux lignes barrées en tête de page, autour du numéro « 55 »}

\emph{Remarque} \quad Inversement, soient $\mathcal{M}$ un ens., $\Phi \subset \mathfrak{P}(\mathcal{M})$ satisf. \add{et} les conditions a) b) c). \emph{NB} La condition c) résulte de a)/b) si $\mathcal{M} \neq \emptyset$. Dans tous les cas, elles équivalent :
\begin{enumerate}
  \item[c')] $\Phi \neq \emptyset$ \quad i.e. existe $A \subset \mathcal{M}$ telle que $A \in \Phi$.
\end{enumerate}
\marginal{\ill{} la condition \ill{} par c \ill{} $A \cap B$ \ill{} (mais $\mathcal{M} \neq \emptyset$ !)}
\note{la note marginale, oblique, en bas à gauche, n'est lue qu'en fragments}

Alors
\begin{enumerate}
  \item[1°)] $\Phi$ ordonné par inclusion satisfait les conditions C 1, C 2. \struck{Ces objets irréductibles}
  \item[2°)] Les objets irréductibles sont les $A_X$, et $X \mapsto A_X$ est une bijection de $\mathcal{M}$ \add{sur} l'ens. des objets irréd. de $\Phi$.
  \item[3°)] Si $X \in \mathcal{M}$, $A \in \Phi$, alors
  \[
    X \in A \Longleftrightarrow A_X \subset A \quad \text{(tautologique)}
  \]
  donc \struck{\ill{}} si on identifie \add{(par $X \mapsto A_X$)} $\mathcal{M}$ à l'ens. des objets irréductibles de $\Phi$, alors pour $A \in \Phi$,
  \[
    \overline{A} = \lbrace A_X \mid A_X \subset A \rbrace \quad \text{i.e. } X \in A
  \]
  s'identifie à $A$ \emph{lui-même}.
\end{enumerate}

On considérera $\mathcal{M}$ comme ordonné par $\leq$ (jamais par $\preccurlyeq$ !) : multistrates et \emph{sous-multistrates}. Les figures \add{$F$} s'identifient aux parties ens. de multistrates $\overline{F}$ correspondantes (cl. par figure à l'ens. de ses strates).

\page{56}
\emph{Proposition} \quad a) Si $F$ \struck{\ill{}} est une figure, alors les \emph{figures} $F'$ telles que $F' \leq F$ s'identifient, par l'application $F' \mapsto \overline{F}'$, aux parties de $\overline{F}$ qui sont \emph{fermées} pour la relation d'ordre \add{$\leq$} de $\mathcal{M}$.

b) Si $F$ et $G$ sont deux figures, $F \wedge G$ existe (inf pour $\leq$) et $\overline{F \wedge G} = \overline{F} \cap \overline{G}$. Même chose pour une famille $(F_i)_{i \in I}$ quelc. \add{avec $I \neq \emptyset$}
\marginal{Mais attention, il n'y a \ill{} $F \wedge G$ \ill{}}
\note{la marge gauche, en haut, porte une note verticale de six lignes, reliée au mot « Proposition », qui n'est lue qu'en fragments ; la mention « avec $I \neq \emptyset$ » est en marge, reliée par un trait}

et \struck{\ill{}} [\emph{Lieux} de $\mathfrak{F}$ : ce sont les \struck{figures} figures qui sont \emph{minimales} dans $\mathfrak{F} \smallsetminus \lbrace \emptyset_{\mathfrak{F}} \rbrace$. \struck{Un lieux est donc} Tout lieu est \add{\uncertain{donc}} irréductible. Donc les lieux sont aussi les éléments minimaux de $\mathcal{M}$ : ce sont les multistrates qui n'ont d'autres strates qu'elle-même.] \struck{Si $F$ \ill{}}
\note{« figures irréductibles » est ajouté en marge, relié au crochet par un trait ; en marge gauche, sur toute la moitié inférieure de la page, une longue note oblique (« \ill{} de $\mathcal{M}$ \ill{} critère \ill{} ») n'est lue qu'en fragments}

\emph{Figures disjointes} $F$ et $G$ : telles que
\begin{enumerate}
  \item[a)] $F$ et $G$ compatibles i.e. $F \cup G$ existe \add{i.e. $\lbrace F, G \rbrace$ majoré pour $\leq$} \struck{et b) on a \ill{}} \struck{(i.e. \ill{} $F \cap G = \emptyset$)}
  \item[b)] $F \cap G = \emptyset_{\mathfrak{F}}$ \struck{\ill{}}
\end{enumerate}

\struck{Si $F$ et $G$ sont}

\emph{Prop} \quad Si $F$ et $G$ sont disjointes, alors pour $F' \leq F$, $G' \leq G$, $F'$ et $G'$ sont disjointes, et $F' \cup G'$ est une sous-figure de $F \cup G$, et l'application $(F', G') \mapsto F' \cup G'$ est une bijection \struck{de l'ens. de $\mathfrak{F}$}

\page{57}
\[
  \underset{\text{ss-figures de } F}{\mathfrak{F}_{\leq F}} \times \underset{\text{ss-fig. de } G}{\mathfrak{F}_{\leq G}}
  \xrightarrow{\ \sim\ }
  \underset{\text{ss-fig. de } F \cup G = H}{\mathfrak{F}_{\leq F \cup G}}
\]
(On récupère $F'$, $G'$ : parties de $H' = F' \cup G'$ comme $H \cap F$ et $H \cap G$.) Itou pour \add{des} \emph{figures} \add{deux à deux disjointes, \ill{}} \struck{\ill{} mutuellement disjointes d'une} configuration (famille pas nécessairement finie). \struck{figure (\ill{})}
\note{on attend $H' \cap F$ et $H' \cap G$ ; la page porte $H$}

\emph{Dévissage} \quad Particulièrement clair, quand on interprète $\mathfrak{F}$ en termes de l'ens. $\Phi \subset \mathfrak{P}(\mathcal{M})$ satisfaisant a) b) c). L'ens. ordonné $\mathcal{M}_{\leq F \vee G}$ s'identifie à la réunion disjointe de l'ens. ordonné somme de $\mathcal{M}_{\leq F}$ et $\mathcal{M}_{\leq G}$, dont ses parties fermées sont \ill{} en \emph{bijection} avec les couples de parties fermées de $\mathcal{M}_{\leq F}$ et $\mathcal{M}_{\leq G}$\ldots

\emph{Remarque} \quad \add{En fin de compte}, la donnée de l'ens. \emph{ordonné} $\mathcal{M}$, \add{et} de $\Phi$ équivalent à : un ens. $\Phi$ de parties \emph{fermées} de $\mathcal{M}$, satisfaisant
\begin{enumerate}
  \item[a)] $\forall A \in \Phi$, \struck{\ill{}} \add{et} partie fermée $A'$ de $A$, on a $A' \in \Phi$
  \item[\struck{b)}] \struck{$\emptyset \in \Phi$ ($\Leftrightarrow \Phi \neq \emptyset$, compte tenu de a)), et}
  \item[b)] $\bigcup_{A \in \Phi} A = \mathcal{M}$ \quad [$\Longleftrightarrow$ \add{compte tenu de a)} $\forall x \in \mathcal{M}$, $\mathcal{M}_{\leq x} \in \Phi$]
  \item[c)] $\emptyset \in \Phi$ \quad ($\Longleftrightarrow$ \add{(a)} $\Phi \neq \emptyset$, et automatiquement satisfait (compte tenu de b)) si $\mathcal{M} \neq \emptyset$.)
\end{enumerate}
\marginal{\ill{} à partir de \ill{} de \ill{}}

\emph{NB} Soit $\mathfrak{P}_f(\mathcal{M})$ l'ens. des parties \emph{fermées} de $\mathcal{M}$, \add{$\mathfrak{P}_f(\mathfrak{F})$ \ill{} de $\mathfrak{F}$ qui sont stables par Sup.} Si $\mathfrak{F}'$ et $\mathcal{M}'$ se correspondent,
\[
  \boxed{\mathfrak{P}_f(\mathfrak{F}) \underset{\psi}{\overset{\varphi}{\rightleftarrows}} \mathfrak{P}_f(\mathcal{M})}
  \qquad
  \varphi(\mathfrak{F}') = \mathfrak{F}' \cap \mathcal{M} , \quad
  \psi(\mathcal{M}') = \lbrace F \in \mathfrak{F} \mid \overline{F} \subset \mathcal{M}' \rbrace
\]
alors $\varphi$, $\psi$ bij. inverses l'une de l'autre. TSVP
\note{ce NB est serré au pied de la page ; la lecture de l'insertion sur $\mathfrak{P}_f(\mathfrak{F})$ est douteuse}

\page{58}
\struck{\emph{C 3} \quad Soient $F$, $G$, $H$ \struck{deux} trois figures mutuellement \add{i.e. deux à deux} compatibles. Alors $\lbrace F, G, H \rbrace$ est une configuration \struck{$F, G, H$} $\exists$ figure $K$ telle que $F$, $G$, $H$ soient des sous-figures de $K$ (donc borne supérieure existe)}
\marginal{alors $\mathfrak{F}'$ satisfait C 1, C 2, \ill{} et ses objets irréductibles \ill{} $\mathcal{M}'$ \ill{} (3) \ill{}}
\note{cette première version de C 3 est barrée de traits obliques ; la note marginale oblique qui l'accompagne, en haut à gauche, n'est lue qu'en fragments}

\emph{C 3} \quad Soient $F, G \in \mathfrak{F}$, telles que $\forall X \in \overline{F}$, $Y \in \overline{G}$, $X$ et $Y$ compatibles. Alors $F$ et $G$ sont compatibles.
\marginal{\emph{Cor} Pour que $F$, $G$ compatibles, il f. et s. que $\forall X \in \overline{F}$, $Y \in \overline{G}$, $X$ et $Y$ le soient}
\note{double trait vertical en marge ; les barres de $\overline{F}$, $\overline{G}$ sont tracées comme de petites flèches}

\emph{NB} L'inverse \uncertain{évident} : si $F$ et $G$ sont compatibles, alors pour toutes sous-figures \add{$F'$ de $F$ (resp.\ ss-fig. id.) et $G'$ de $G$}, $F'$ et $G'$ le sont.
\note{l'insertion est écrite au-dessus de la ligne et reprise en tête de la suivante ; sa construction est celle de la page}

\emph{Cor} \quad Soient $F$ et $G$, $H$ \struck{deux} \add{trois} figures, \struck{\ill{}} $F$, $G$ compatibles. Pour que $H$ soit compatible à $F \cup G$, il f. et s. qu'il le soit à $F$ et à $G$. Plus gén., soit $(F_i)_{i \in I}$ une configuration, $F$ la figure \struck{\ill{}} réunion. Pour que $H$ soit compatible à $F$, il f. et s. que $H$ soit compatible aux $F_i$.
\note{on lit $F_i$ à la dernière ligne, qui porte $F$ suivi d'un signe}

\page{59}
Donc, en termes d'une structure portée par un ens. ordonné $\mathcal{M}$ (cf. remarque page 58), on trouve ceci. On a sur \struck{Définition} $\mathcal{M}$ une relation $R$ (« compatibilité ») qui est réflexive symétrique (donc \uncertain{équivalente} \ill{} et) telle que \ill{} $S \subset \mathfrak{P}_2(\mathcal{M})$), \struck{implique} la donnée de $R(X, Y)$ implique $R(X', Y')$ pour $X' \leq X$, $Y' \leq Y$. \struck{compatible avec $\leq$}. Cela permet, à peu de choses près, de reconstituer $(\mathfrak{F}, \leq)$, en tant que ens. de parties de $\mathcal{M}$. Pour l'explication,
\note{« page 58 » renvoie sans doute à la remarque de la page 57, qui s'achève sur « TSVP » ; la parenthèse sur $S \subset \mathfrak{P}_2(\mathcal{M})$ est écrite au-dessus de la ligne et n'est lue qu'en partie}

Quelques terminologies (\add{strates maximales}) « \emph{Composantes irréductibles} » d'une figure $F$ : strates maximales de $F$ ($\exists$ toujours si l'ens. $\overline{F}$ des strates de $F$ est fini (on dit alors que $F$ est une \emph{figure finie}) et $F \neq \emptyset_{\mathfrak{F}}$).

Figure \struck{de type} \add{de type} fini : qui est borne sup. de ses « composantes » irréductibles en nombre fini (\uncertain{qui} est de t.f.), (exemple : toute figure finie $F$ / qui est \struck{\ill{}} borne sup. \struck{\ill{} d'une \ill{} $\operatorname{Sup} X$ \ill{}} \ill{} figures irréductibles. \ill{} (\emph{NB} Dans les figures les plus importantes, toute figure est finie, et a fortiori de type fini.)
\note{la phrase entre « exemple » et « irréductibles » est corrigée et en partie barrée ; elle n'est lue qu'en partie}

Ceci dit si $\mathfrak{F}_{\mathrm{t.f.}}$ désigne l'ens. \emph{figures} $\mathfrak{F}_{\mathrm{t.f.}}$, l'ens. de type fini, on récupère \ill{} ens. de parties correspondant de $\mathcal{M}$, par la seule connaissance de $(\mathcal{M}, \leq, R)$,

\page{60}
de la façon suivante : ce sont les $A \subset \mathcal{M}$ tels que $A = \bigcup_i A_{X_i}$, où $\lbrace X_i \rbrace$ \struck{\ill{}} est un ens. fini d'éléments de $\mathcal{M}$, mutuellement compatibles (au sens de $R$).
\note{la page porte $A \in \mathcal{M}$}

Donc une structure $(\mathfrak{F}, \leq)$ satisfaisant C 1, C 2, C 3 + figures toutes de type fini, équivaut à $(\mathcal{M}, \leq, R)$, où $R$ rel. symétrique réfl. \add{i.e. reliée à $\leq$ et $R$} satisfaisant \struck{deux} conditions :
\begin{enumerate}
  \item[a)] \struck{\ill{}} (\uncertain{condition} \ill{} $\leq$) Pour $A$ partie fermée \emph{majorée} $A$ de $\mathcal{M}$, $A$ est de type fini, i.e. de la forme $\bigcup \mathcal{M}_{\leq X_i}$, où les $X_i \in A$ forment un ens. fini [\ill{} il suffit que les $\mathcal{M}_{\leq X}$ soient finis]
  \item[b)] (compatibilité de $\leq$ \emph{et} $R$) si $(X, Y) \in R$, et $X' \leq X$, $Y' \leq Y$, alors $(X', Y') \in R$.
\end{enumerate}
\marginal{\ill{} donc \ill{} fermée \ill{} fini}
\note{la note marginale oblique, en face de a), n'est lue qu'en fragments ; le crochet final de a) est écrit en marge droite}

Les « figures » de $\mathcal{M}$ sont alors les parties \struck{fermées} \add{$A$} de $\mathcal{M}$ qui ont les propriétés suivantes :
\begin{enumerate}
  \item[a)] $A$ est fermée
  \item[b)] $A$ est de type fini
  \item[c)] $\forall X, Y \in A$, on a $(X, Y) \in R$.
\end{enumerate}
\marginal{(\ill{} propriété Cont 2 de p. 33)}
\note{« p. 33 » renvoie à la page 33, au deuxième lot}
\end{document}
